\subsection{Linear mappings of free modules, and matrices over a principal ideal domain}

Let A be a principal ideal domain. Consider a linear mapping f from a free A-module L of rank m into a free A-module \(L'\) of rank n.~The preceding results allow us, by choosing suitable bases for L and \(L'\) , to put the matrix for f in a particularly simple form, called the canonical form of the matrix.

PROPOSITION 5. --- Let A be a principal ideal domain, and f a linear map of rank r from a free A-module L of rank m into a free A-module L' of rank n.~Then there exist bases \((e_{i})\) \((1 \leqslant i \leqslant m)\) of L and \((e_{j}')\) \((1 \leqslant j \leqslant n)\) of L' such that \(f(e_{i}) = \alpha_{i} e_{i}'\) for \(1 \leqslant i \leqslant r\) and \(f(e_{i}) = 0\) for i \textgreater{} r, where the \(\alpha_{i}\) are nonzero elements of A, each of which divides the next; the ideals \(A\alpha_{i}\) are the invariant factors of \(f(L)\) in \(L'\) , and are thus uniquely determined.

Let \(L_0 = f^{-1}(0)\) be the kernel of \(f\) ; the quotient \(L/L_0\) is isomorphic to the module \(f(L)\) , which is a submodule of \(L'\) and so free (VII, p.~15, Cor. 2); thus \(L_0\) has a complement \(L_1\) in \(L\) (II, p.~218, Prop. 21), and the restriction of \(f\) to \(L_1\) is an isomorphism from \(L_1\) onto \(f(L) = M'\) . If the ideals \(A\alpha_i\) ( \(1 \leqslant i \leqslant r\) ) are the invariant factors of \(M'\) in \(L'\) , then Th. 1 of VII, p.~18 shows that there exists a basis \((e_j')\) ( \(1 \leqslant j \leqslant n\) ) of \(L'\) such that \((\alpha_i e_i')\) ( \(1 \leqslant i \leqslant r\) ) is a basis of \(M'\) . Since the restriction of \(f\) to \(L_1\) is an isomorphism from \(L_1\) onto \(M'\) , there exists a basis \((e_i)\) ( \(1 \leqslant i \leqslant r\) ) of \(L_1\) such that \(f(e_i) = \alpha_i e_i'\) . This basis extends to a basis \((e_k)\) ( \(1 \leqslant k \leqslant m\) ) of \(L\) by taking \((e_s)\) ( \(r + 1 \leqslant s \leqslant m\) ) to be a basis of the kernel

\[
\mathbf {L} _ {0}
\]

COROLLARY 1. --- Let X be a matrix of rank r, with n rows and m columns, over a principal ideal domain A; then there exists a matrix \(X_{0}\) equivalent to X (II, p.~354) of the form

\[
\left( \begin{array}{c c c c c c c} \alpha_ {1} & 0 & \ldots & 0 & 0 & \ldots & 0 \\ 0 & \alpha_ {2} & \ldots & 0 & 0 & \ldots & 0 \\ 0 & 0 & \ldots & \alpha_ {r} & 0 & \ldots & 0 \\ 0 & 0 & \ldots & 0 & 0 & \ldots & 0 \\ 0 & 0 & \ldots & 0 & 0 & \ldots & 0 \end{array} \right)
\]

where the \(\alpha_{i}\) are nonzero elements of A, each of which divides the next. Under these conditions the \(\alpha_{i}\) are uniquely determined up to multiplication by invertible elements.

Given that two matrices X and \(X'\) are equivalent if there exist invertible square matrices P and Q, of orders n and m, over A, such that \(X' = PXQ\) , Cor. 1 is just Prop. 5 expressed in terms of matrices.

In the notation of Prop. 5 and Cor. 1, the (nonzero) ideals \(\mathbf{A}\alpha_{i}\) are called the invariant factors of the linear map \(f\) , or of the matrix \(X\) . It then follows immediately from Cor. 1 that:

COROLLARY 2. --- Two matrices X and \(X'\) with n rows and m columns over a principal ideal domain A are equivalent if and only if they have the same invariant factors.

Note that when A is a field we can take the \(\alpha_{i}\) to be equal to 1, and then we recover Prop. 13 of II, p.~360.

If X is the matrix of the linear mapping f with respect to an arbitrary basis of L and an arbitrary basis of \(L'\) , then the columns of X are the coordinate vectors, with respect to the basis of \(L'\) , of elements of \(L'\) which form a system of generators for \(f(\mathrm{L})\) . The following result is thus an immediate consequence of Prop. 4.

PROPOSITION 6. --- Let X be a matrix of rank r over a principal ideal domain A, and let \(A\alpha_{i}\) ( \(1 \leqslant i \leqslant r\) ) be its sequence of invariant factors. Then \(\alpha_{1}\) is a gcd of the elements of X; and the product \(\alpha_{1} \ldots \alpha_{q}\) is a gcd of the q-th order minors of X for \(q \leqslant r\) .

